\section{Pretraining, Cost, and External Evaluation}
\label{app:systems}

\subsection{Measured Policy Cost and Physical Delay}
The future-screen policy call means span 177.071--184.212\,ms;
p95 spans 178.726--185.133\,ms; peak allocated memory spans
1.571--1.617\,GiB. These are same-device, fixed-order measurements
including internal perception and action computation. Loading,
external oracle reads, I/O, rendering, and priming are excluded.
Each arm uses three recorded eight-frame sequences, six warm-up calls,
and thirty measured calls (ten per sequence). After each reset, seven real
calls prime the observation history; the eighth is timed with CUDA synchronization
before and after the call. The sequences cover translation, occlusion,
and post-impact replanning. Inference is FP32 with TF32 enabled,
batch/sequence dimensions $B=S=1$, and five Euler action steps.
NVML identifies the shared device as \texttt{NVIDIA RPBZZZ6}, with
102,641,958,912 bytes of memory; CPU affinity is 0--2.
Memory is the post-warm-up PyTorch allocation, including resident weights.
One measurement round supports this range;
differences of a few milliseconds do not establish a robust speed ranking.

For deployed commands, action age is the difference between sensor capture
and command application. The planned deployment profile separates capture,
availability, preprocessing, inference, queuing, issue, and application.
It records hardware, precision, batch size, warm-up, synchronization,
and timed sample counts. Fixed-delay simulation tests a controlled lag;
wall-clock execution tests the actual combined system.
The final same-device profiles, deadline misses, and action-age distributions
are \placeholder{pending}.

\subsection{Recipe and Robot-Data Base}
After the design cross-checks, the study fixes input modalities, history
support, future targets/horizons, information flow, and action semantics.
The next stage optimizes the existing loss weights and training recipe,
and changes depth/width within the chosen architecture for scale comparison.
Changing the representation or adding an information path belongs to a
new design comparison rather than an optimizer ablation.

DROID supplies real robot demonstrations \cite{droid}. Preparation checks
sensor timing, calibration, depth, command versus measured motion,
gripper semantics, and training-only normalization. It does not provide
object pose GT; the target-pose pipeline must be separately validated
for any selected pose-based objective. Only targets used by the selected
design and necessary ablations need to be produced.
The data report includes valid-label masks and rejected-action rates.
This stage is \placeholder{planned}.

The formal base is trained on real robot data. Two independent post-training
branches start from the same base: one for simulation and one for
real-robot tasks. Simulation-from-scratch experiments supply recipe
experience and a strong reference, while their weights do not initialize
the formal robot-data base or the real-robot branch.
Comparisons with and without pretraining share downstream data and
exposure, and additionally allow equally budgeted recipe searches.
They report the pretraining cost explicitly. Data, model, and compute
scale have separate axes.

\subsection{Planned External and Real-Robot Evidence}
Single-arm DOMINO \cite{domino} tests dynamics under a different task/data
protocol. Static compatibility uses an explicitly selected native suite
such as LIBERO \cite{libero}. Real conveyor manipulation evaluates
perception, prediction, action, and timing together.
The final phase reuses completed pretraining, scaling, and independent
adaptation results, then fills uncovered object, scene, motion, and delay
conditions. It does not repeat an already completed scaling matrix.

The real-robot protocol must specify objects, speed/delay strata, trial
allocation, ordering, calibration, safety termination, and human intervention.
External native scores are reported separately from the fourteen-task
simulation SR. Final training scale, robot trials, and real-world performance
are \placeholder{pending}.
